Assume also that  $x$  has an entry which is a multiple of $\ell$. 
Choose a minimal entry of $x$  
which is a multiple of $\ell$. 
If there are several such entries,  
take one with the largest index. 
This index  $i \in [n] = \{1,\ldots,n\}$  
will be called the {\em pivot} 
and $x_i$ the {\em pivotal entry} of $x$. 
% Obviously, the pivot is unique.  

\medskip 

The {\em GM-rule} is defined as follows: 
keep the pivotal entry unchanged 
and reduce the remaining $n-1$ entries by 1.
This transformation is called the {\em GM-move} 
\subsection{Game NIM$(n,n-1)$  for $\ell$ players} 
The exact slow NIM$(n,k)$, with $k = n-1$  
can be played by $\ell$  players  % $j_1, \ldots, j_\ell$  
as follows. 
A position is a non-negative $n$-vector  $x$. 
% $x = (x_1, \ldots, x_n)$. 
Players make moves in a given cyclic order. 
By one move a player can choose an arbirary 
entry and keep it unchanged 
reducicing $n-1$ remaining entries by 1, 
provided they are positive. 
If  $x$  has at least two non-positive entries, 
it is called a {\em terminal} position, 
the game is over, and the player 
who has to move in $x$ (but cannot) lost, 
while  $\ell-1$  others won. 
The payoffs are defined below. 
A sequence of moves 
from the initial position  $x^0$ 
to a terminal one is called a {\em play}. 
Let  $L(P)$  denote the length, 
that is, the number of moves, of a play  $P$. 
Choose a large constant $C$; 
it should be larger than the length 
of any play from  $x^0$.  
Then the loser pays $C-L(P)$ 
and $\ell-1$  winners share this amount, 
genning $\frac{C - L(P)}{\ell-1}$ each. 
The GM-rule defines the unique strategy foreach player.
If a position  $x$  has no entries which are multiple of $\ell$, 
the GM-move in $x$ keeps the lagest entry and reduces 
the smaller $n-1$ by 1 each. 

\begin{conjecture}
\label{con-1} 
The set of $\ell$ GM-strategies form a (uniform) Nash equilibrium. 
\end{conjecture}

By definition, the GM-strategies are uniform, 
that is, independent of  $x^0$. 
Hence, a  Nash equilibrium in GM-strategies  
is uniform, if exists. 

